\documentclass[aspectratio=169]{beamer}
\usetheme{Madrid}
\usecolortheme{whale}

\usepackage[utf8]{inputenc}
\usepackage{graphicx}
\usepackage{booktabs}
\usepackage{listings}
\usepackage{amsmath}
\usepackage{amssymb}
\usepackage{hyperref}
\usepackage{color}

% Custom Premium Indigo/Slate Color Theme (Professional & Sleek)
\definecolor{slate}{RGB}{30, 41, 59}
\definecolor{indigo}{RGB}{79, 70, 229}
\definecolor{teal}{RGB}{13, 148, 136}

\setbeamercolor{palette primary}{bg=slate,fg=white}
\setbeamercolor{palette secondary}{bg=indigo,fg=white}
\setbeamercolor{palette tertiary}{bg=slate,fg=white}
\setbeamercolor{titlelike}{parent=palette primary}
\setbeamercolor{structure}{fg=indigo}

\lstset{
  basicstyle=\ttfamily\tiny,
  keywordstyle=\color{indigo}\bfseries,
  commentstyle=\color{gray},
  stringstyle=\color{teal},
  showstringspaces=false,
  breaklines=true
}

% Presentation Details
\title[Full-Stack Geospatial Engineering]{Full-Stack Geospatial Engineering \& Data Pipelines}
\subtitle{Scaling Decoupled Media Workspaces for Aerial Infrastructure Inspections}
\author{Rares Olteanu}
\date{June 2026}

\begin{document}

% Slide 1: Title Slide
\begin{frame}
  \titlepage
\end{frame}

% Slide 2: Agenda
\begin{frame}{Agenda}
  \begin{columns}
    \begin{column}{0.5\textwidth}
      \textbf{Part 1: Ingestion & Frontend (Slides 4-12)}
      \begin{itemize}
        \item EXIF & Camera Telemetry Extraction
        \item Debounced Viewport Queries
        \item Leaflet Marker Clustering
      \end{itemize}
      \textbf{Part 2: Backend & AI Pipelines (Slides 13-21)}
      \begin{itemize}
        \item Express API & Relational Indexes
        \item Asynchronous Background Workers
        \item Vision AI Error Resilience
      \end{itemize}
    \end{column}
    \begin{column}{0.5\textwidth}
      \textbf{Part 3: Security & Code Quality (Slides 22-26)}
      \begin{itemize}
        \item Cryptographic CSRNG Verification
        \item Dead Code & Refactoring Actions
      \end{itemize}
      \textbf{Part 4: Testing & Pipelines (Slides 27-31)}
      \begin{itemize}
        \item Mocking Network Latencies with MSW
        \item Dockerization & CI/CD Workflows
      \end{itemize}
      \textbf{Part 5: Production Scaling (Slides 32-40)}
      \begin{itemize}
        \item PostgreSQL/PostGIS Spatial Indexes
        \item Bounding Boxes & Proximity Queries
      \end{itemize}
    \end{column}
  \end{columns}
\end{frame}

% Slide 3: Candidate Profile & Philosophy
\begin{frame}{Candidate Profile \& Philosophy}
  \begin{itemize}
    \item \textbf{Core Expertise:} Full-Stack Software Engineering, API Design, and Geospatial Data Pipelines.
    \item \textbf{Software Quality Advocacy:} Focus on writing clean, self-documenting code, enforcing static analysis, and eliminating dead code paths.
    \item \textbf{Robust Testing Standards:} Integrating MSW network mocking, Unit testing, and CI/CD pipelines to ensure release reliability.
    \item \textbf{System Vision:} Designing local prototypes with a clear, documented path to cloud-scale, highly available architectures.
  \end{itemize}
\end{frame}

% Slide 4: Ingestion Context: Mapping Aerial Inspection Media
\begin{frame}{Mapping Aerial Inspection Media}
  \begin{columns}
    \begin{column}{0.6\textwidth}
      \begin{itemize}
        \item \textbf{The Challenge:} Inspecting linear critical infrastructures (pipelines, power lines, rail grids) generates thousands of high-resolution images.
        \item \textbf{Target Capabilities Needed:}
          \begin{itemize}
            \item Seamless metadata mapping to geographic coordinates.
            \item High-speed visualization of media points without UI freezing.
            \item Context-aware description and anomaly tags.
            \item Decoupled, non-blocking upload cycles.
          \end{itemize}
      \end{itemize}
    \end{column}
    \begin{column}{0.4\textwidth}
      \begin{center}
        \textbf{Visual Imagery} \\
        $\downarrow$ \\
        \textbf{EXIF Coordinate Extraction} \\
        $\downarrow$ \\
        \textbf{Interactive Geospatial Map}
      \end{center}
    \end{column}
  \end{columns}
\end{frame}

% Slide 5: Decoupled Product Workflow
\begin{frame}{Decoupled Product Workflow}
  \begin{itemize}
    \item \textbf{Dynamic Frontend Client:} A React SPA that handles heavy EXIF reading, local coordinate overrides, and Leaflet map rendering.
    \item \textbf{Express API Server:} Lightweight REST server handling auth verification, database queries, and media formatting.
    \item \textbf{Asynchronous Worker Processes:} Database coordinates are stored instantly, while heavy Vision AI jobs are enqueued in the background.
    \item \textbf{Decoupled Assets:} Serve compressed images and thumbnails directly to users while preserving original files.
  \end{itemize}
\end{frame}

% Slide 6: Frontend Tech Stack
\begin{frame}{Frontend Tech Stack}
  \begin{table}
    \centering
    \begin{tabular}{lll}
      \toprule
      \textbf{Layer} & \textbf{Technology} & \textbf{Core Purpose} \\
      \midrule
      Runtime Framework & React / Vite & Ultra-fast HMR \& Component-Driven UI \\
      Geospatial Rendering & Leaflet / React Leaflet & Interactive maps with custom photo elements \\
      Clustering Engine & React Leaflet Cluster & Merges dense nodes on zoom shifts \\
      Metadata Extractor & exifr & Parses coordinate data on client thread \\
      HTTP client & Axios & Reusable requests with JWT authorization \\
      Testing Harness & Vitest / React Testing Library & Component-level verification \\
      \bottomrule
    \end{tabular}
  \end{table}
\end{frame}

% Slide 7: Frontend Component Architecture
\begin{frame}{Frontend Component Architecture}
  \begin{itemize}
    \item \textbf{MapPage (Dashboard):} Stores global layout, search keywords, selected media states, and visible viewport coordinate bounds.
    \item \textbf{MapView (Geospatial Panel):} Embeds the Leaflet container, dynamic marker cluster layers, and viewport movement event listeners.
    \item \textbf{UploadModal (Ingestion Panel):} Handles drag-and-drop file inputs, calls the EXIF parser, and displays coordinate fallbacks.
    \item \textbf{PhotoModal (Detail Panel):} Renders large image previews, displays AI descriptions, hosts comments threads, and handles owner deletion actions.
    \item \textbf{AuthPage (Access Control):} Houses JWT token inputs, signup verify screens, and password reset flows.
  \end{itemize}
\end{frame}

% Slide 8: Client-Side Telemetry Extraction: EXIF Coordinate Parsing
\begin{frame}{Client-Side Ingestion: EXIF Coordinate Parsing}
  \begin{itemize}
    \item \textbf{Concept:} Photos captured by GPS-enabled sensors contain rich metadata in EXIF headers (Exchangeable Image File Format).
    \item \textbf{Client-First Ingestion:}
      \begin{itemize}
        \item Frontend reads the file locally using `exifr` before upload.
        \item Parses latitude and longitude tags, converting degrees-minutes-seconds to decimal formats.
        \item Autofills latitude and longitude coordinates in the upload form.
      \end{itemize}
    \item \textbf{Performance Benefit:} Offloads metadata parsing overhead from the API server and gives immediate visual coordinate feedback to the user.
  \end{itemize}
\end{frame}

% Slide 9: Advanced EXIF: Reading Camera Yaw/Pitch/Roll & Altitude
\begin{frame}{Advanced EXIF: Reading Camera Yaw/Pitch/Roll \& Altitude}
  \begin{itemize}
    \item \textbf{Inspection Requirements:} In infrastructure surveys (e.g. airships), knowing only the GPS position of the sensor is not enough. We also need sensor orientation.
    \item \textbf{Telemetry Extraction:}
      \begin{itemize}
        \item \textbf{Altitude:} Identifies the height of the sensor (crucial for perspective calculation).
        \item \textbf{Gimbal Angles (Yaw, Pitch, Roll):} Stored in XMP metadata schemas. Denotes the exact direction the camera was pointing relative to north and the horizon.
      \end{itemize}
    \item \textbf{Application:} Allows projecting the exact footprint of the photo onto the map, rather than rendering a simple point marker.
  \end{itemize}
\end{frame}

% Slide 10: Smart User Fallbacks: Interactive Mini-Map Marker Placement
\begin{frame}{Smart User Fallbacks: Interactive Mini-Map}
  \begin{itemize}
    \item \textbf{The Problem:} Over $50\%$ of web images have GPS metadata stripped due to privacy rules or sensor misconfigurations.
    \item \textbf{UX Solution:}
      \begin{itemize}
        \item If EXIF parsing fails to return coordinates, the upload modal displays an interactive fallback interface.
        \item Users can type decimal coordinates directly or click on an embedded mini Leaflet map.
        \item Map clicks fire coordinate events, placing a visual marker and updating form inputs instantly.
      \end{itemize}
    \item \textbf{Benefit:} Guarantees that every uploaded photo remains associated with a valid spatial point.
  \end{itemize}
\end{frame}

% Slide 11: Interactive Map Rendering: Leaflet Marker Clustering
\begin{frame}{Interactive Map Rendering: Leaflet Marker Clustering}
  \begin{itemize}
    \item \textbf{The Bottleneck:} Rendering hundreds of individual DOM markers in close proximity crashes browser rendering threads.
    \item \textbf{Marker Clustering Strategy:}
      \begin{itemize}
        \item Integrates `react-leaflet-cluster` (utilizing the high-performance Supercluster library).
        \item Groups close markers into a single cluster node containing a count badge.
        \item Zooming in dynamically expands clusters into sub-clusters or individual photo markers.
      \end{itemize}
    \item \textbf{Benefit:} Maintains high frame rates (60 FPS) when navigating dense geospatial regions.
  \end{itemize}
\end{frame}

% Slide 12: Client-Side Bounding Box Event Handlers
\begin{frame}{Client-Side Bounding Box Event Handlers}
  \begin{itemize}
    \item \textbf{Core Concept:} The map should not load the entire global photo database. It must query only what is currently visible on the screen.
    \item \textbf{Event-Driven Fetching:}
      \begin{itemize}
        \item Map fires Leaflet `moveend` and `zoomend` events.
        \item The frontend captures the active map bounding box:
          \[ \text{bbox} = \{\text{South}, \text{West}, \text{North}, \text{East}\} \]
        \item Passes coordinates to the API via `GET /photos?bbox=south,west,north,east`.
      \end{itemize}
    \item \textbf{Optimization:} The query is debounced to prevent duplicate network calls during rapid zooming or panning.
  \end{itemize}
\end{frame}

% Slide 13: Backend Tech Stack
\begin{frame}{Backend Tech Stack}
  \begin{table}
    \centering
    \begin{tabular}{lll}
      \toprule
      \textbf{Layer} & \textbf{Technology} & \textbf{Core Purpose} \\
      \midrule
      API Server & Node.js / Express & Lightweight, high-throughput REST API \\
      Database Engine & SQLite / LibSQL & Local, low-latency relational database \\
      Media Processing & Sharp & Image resizing, rotation, and metadata standardization \\
      Authentication & jsonwebtoken / bcrypt & Secure JWT signing & password hashing \\
      Mailer & Nodemailer & Dispatches verification codes (Ethereal Sandbox) \\
      AI Integration & @google/generative-ai & Offloaded Vision AI content generation \\
      \bottomrule
    \end{tabular}
  \end{table}
\end{frame}

% Slide 14: REST API Endpoint Design
\begin{frame}{REST API Endpoint Design}
  \begin{itemize}
    \item \textbf{Decoupled REST Routing:} Organized logically around resource patterns.
    \item \textbf{Public Endpoints:}
      \begin{itemize}
        \item `POST /auth/send-signup-code` / `POST /auth/signup` (Registration flow)
        \item `POST /auth/login` (Returns JWT access token)
      \end{itemize}
    \item \textbf{JWT-Protected Endpoints (enforcing authMiddleware):}
      \begin{itemize}
        \item `GET /photos` (Supports optional `bbox` coordinate parameters)
        \item `POST /photos` (Handles multipart form uploads)
        \item `DELETE /photos/:id` (Checks owner credentials before deleting)
        \item `POST /photos/:id/regenerate-description` (Triggers AI updates)
        \item `POST /comments/:photoId` (Adds user comment to thread)
      \end{itemize}
  \end{itemize}
\end{frame}

% Slide 15: Relational Data Model
\begin{frame}{Relational Data Model}
  \begin{itemize}
    \item \textbf{Schema Structure:} Designed as a clean relational model to maintain referential integrity.
    \item \textbf{Core Relations:}
      \begin{itemize}
        \item \textbf{users:} Unique ID, email (unique index), password hash, display name, and username.
        \item \textbf{photos:} Linked to `users` (FK). Stores filename, original metadata name, decimal coordinates, and AI description string.
        \item \textbf{comments:} Linked to `photos` (FK) and `users` (FK). Stores text content and creation timestamp.
      \end{itemize}
    \item \textbf{Cascade Deletions:} Deleting a photo automatically cascades and cleans up associated comments in the database.
  \end{itemize}
\end{frame}

% Slide 16: Relational Indexing Strategy for Fast Lookups
\begin{frame}{Relational Indexing Strategy for Fast Lookups}
  \begin{itemize}
    \item \textbf{Query Performance:} As databases grow, scanning full tables for spatial bounding boxes or user sessions becomes a major bottleneck.
    \item \textbf{Index Allocations:}
      \begin{itemize}
        \item `CREATE UNIQUE INDEX idx_users_email ON users(email)` $\to$ Speeds up authentication checks.
        \item `CREATE INDEX idx_photos_coords ON photos(lat, lng)` $\to$ Optimizes bounding box calculations.
        \item `CREATE INDEX idx_photos_user ON photos(user_id)` $\to$ Speeds up user gallery views.
        \item `CREATE INDEX idx_comments_photo ON comments(photo_id, created_at)` $\to$ Renders comment threads instantly.
      \end{itemize}
  \end{itemize}
\end{frame}

% Slide 17: Media Pipeline: Normalization & Compression using Sharp
\begin{frame}{Media Pipeline: Normalization \& Compression}
  \begin{itemize}
    \item \textbf{The Orientation Bug:} Mobile photos contain EXIF orientation markers. Standard web browsers often render these rotated.
    \item \textbf{Sharp Integration:}
      \begin{itemize}
        \item The upload router intercepts file streams using `multer`.
        \item Passes image buffers to `sharp`.
        \item Executes `.rotate()` to read EXIF orientation metadata and permanently rotate the pixel grid.
        \item Converts images to JPEG format with $80\%$ quality compression.
      \end{itemize}
    \item \textbf{Result:} Normalizes all uploads for uniform browser display while reducing storage consumption by up to $60\%$.
  \end{itemize}
\end{frame}

% Slide 18: Asynchronous Processing Pipeline
\begin{frame}{Asynchronous Processing Pipeline}
  \begin{itemize}
    \item \textbf{The Design Challenge:} Requesting Vision AI descriptions during upload blocks the HTTP response cycle, forcing users to wait up to 10 seconds.
    \item \textbf{Asynchronous Worker Architecture:}
  \end{itemize}
  \begin{center}
    Upload File $\to$ Save DB Entry $\to$ Send HTTP 201 Response $\to$ Spawn Background AI Generator
  \end{center}
  \begin{itemize}
    \item \textbf{Advantage:} Users get immediate feedback. The frontend renders the new photo marker instantly, and the description fills in asynchronously.
  \end{itemize}
\end{frame}

% Slide 19: Asynchronous AI Architecture: Vision Model Integration
\begin{frame}[fragile]{Asynchronous AI Architecture: Vision Model Integration}
  \begin{itemize}
    \item \textbf{AI Engine:} Uses the multimodal Google Gemini API to analyze visual features.
    \item \textbf{Implementation Details:}
  \end{itemize}
\begin{lstlisting}[language=JavaScript]
async function generateAIDescription(filePath) {
  if (!process.env.GEMINI_API_KEY) return null;
  const model = genAI.getGenerativeModel({ model: "gemini-1.5-flash" });
  const result = await model.generateContent([
    "Provide a 1-sentence descriptive caption of this photo.",
    { inlineData: { data: fs.readFileSync(filePath).toString("base64"), mimeType: "image/jpeg" } }
  ]);
  return result.response.text();
}
\end{lstlisting}
\end{frame}

% Slide 20: AI Resilience: Quota Handling & Exponential Backoff
\begin{frame}{AI Resilience: Quota Handling \& Exponential Backoff}
  \begin{itemize}
    \item \textbf{Production Reality:} Third-party APIs face transient network drops, rate limits (429), or server congestion (503).
    \item \textbf{Retry Strategy with Exponential Backoff:}
      \begin{itemize}
        \item If the Gemini API call fails, the helper catches the exception and waits.
        \item Retry limit set to 3.
        \item Delays double on each attempt: $2\text{s} \to 4\text{s} \to 8\text{s}$.
      \end{itemize}
    \item \textbf{Graceful Degradation:} If all retries fail, the description defaults to `NULL` in the database, ensuring the core upload workflow is unaffected.
  \end{itemize}
\end{frame}

% Slide 21: Dynamic Regeneration: On-Demand AI Updates
\begin{frame}{Dynamic Regeneration: On-Demand AI Updates}
  \begin{itemize}
    \item \textbf{Use Case:} If AI fails during initial upload (or if API keys were not configured yet), users shouldn't have to re-upload.
    \item \textbf{Interactive Trigger:}
      \begin{itemize}
        \item The photo detail modal displays a "Regenerate Description" button.
        \item Invokes `POST /photos/:id/regenerate-description`.
        \item Runs the AI pipeline, updates the database row, and returns the description to update the frontend state.
      \end{itemize}
    \item \textbf{Benefit:} Gives users control and provides a fallback mechanism for temporary API outages.
  \end{itemize}
\end{frame}

% Slide 22: Stateless Authentication: JWT Verification Middleware
\begin{frame}{Stateless Authentication: JWT Verification}
  \begin{itemize}
    \item \textbf{Concept:} Session states are not stored on the server. The server verifies incoming requests mathematically.
    \item \textbf{JWT Flow:}
      \begin{itemize}
        \item Successful login returns a token signed using a server-side secret (`JWT_SECRET`).
        \item Token contains user ID and email credentials.
      \end{itemize}
    \item \textbf{Guard Middleware:}
      \begin{itemize}
        \item Intercepts requests, parses the `Authorization: Bearer <token>` header.
        \item Calls `jwt.verify()`.
        \item If valid, extracts payload and passes control to the route handler. Otherwise, returns a `401 Unauthorized` response.
      \end{itemize}
  \end{itemize}
\end{frame}

% Slide 23: Cryptographically Secure OTP Generation
\begin{frame}{Cryptographically Secure OTP Generation}
  \begin{itemize}
    \item \textbf{Vulnerability Fixed:} Standard random functions (`Math.random()`) are pseudo-random and vulnerable to prediction attacks.
    \item \textbf{Implementation Upgrade:}
      \begin{itemize}
        \item Replaced Math.random with Node's native `crypto.randomInt()`.
        \item Generates high-entropy 8-digit verification codes.
      \end{itemize}
    \item \textbf{Why it matters:} Ensures that verification and reset tokens cannot be guessed, establishing a secure baseline for registration and password resets.
  \end{itemize}
\end{frame}

% Slide 24: Code Lifecycle Security: TTL & Single-Use Cleanup
\begin{frame}{Code Lifecycle Security: TTL \& Single-Use Cleanup}
  \begin{itemize}
    \item \textbf{State Control:} Codes are stored in an in-memory mapping with security constraints.
    \item \textbf{Expiry Time-to-Live (TTL):}
      \begin{itemize}
        \item Active lifetime is restricted to 10 minutes (`expiresAt: Date.now() + 10 * 60 * 1000`).
        \item Expired codes are rejected and removed from memory.
      \end{itemize}
    \item \textbf{Single-Use Deletion:}
      \begin{itemize}
        \item Immediately deletes the stored code record upon first successful validation.
      \end{itemize}
    \item \textbf{Benefit:} Prevents brute-force attempts and replay attacks.
  \end{itemize}
\end{frame}

% Slide 25: Code Quality Refactoring: Clean Architecture
\begin{frame}{Code Quality Refactoring: Clean Architecture}
  \begin{itemize}
    \item \textbf{Refactoring Operations:}
      \begin{itemize}
        \item Cleaned frontend state logic by removing dead variables (`verificationCode`) and client-side code generation.
        \item Structured backend folders: `routes/` (endpoints), `middleware/` (guards), `utils/` (email helpers).
        \item Removed informal debug strings and signatures from production mailers.
      \end{itemize}
    \item \textbf{Results:} Avoids compiler and linter warnings, simplifies codebase onboarding, and presents a professional standard to reviewers.
  \end{itemize}
\end{frame}

% Slide 26: Environment Portability: Configurable APIs
\begin{frame}{Environment Portability: Configurable APIs}
  \begin{itemize}
    \item \textbf{The Issue:} Hardcoded environment variables require code updates for different deployments.
    \item \textbf{Solution:} Moved all configurations to environment variables:
      \begin{itemize}
        \item `VITE_API_BASE_URL` dynamically routes the frontend client.
        \item `FRONTEND_URL` configures CORS origins on the backend.
        \item `GEMINI_API_KEY` enables or disables the AI subsystem.
        \item `SMTP_HOST` / `SMTP_USER` / `SMTP_PASS` manages email routing (with automatic fallback to Ethereal sandbox accounts).
      \end{itemize}
    \item \textbf{Benefit:} Code remains unchanged. Environment-specific settings are injected at runtime.
  \end{itemize}
\end{frame}

% Slide 27: Testing Philosophy: Unit vs. E2E Tests
\begin{frame}{Testing Philosophy: Unit vs. E2E Tests}
  \begin{itemize}
    \item \textbf{Scope:} Testing is split into distinct strategies to ensure code stability.
    \item \textbf{Unit Tests (Vitest + React Testing Library):}
      \begin{itemize}
        \item Fast execution (runs in milliseconds).
        \item Verifies component-level logic (e.g. upload coordinates update when EXIF succeeds).
      \end{itemize}
    \item \textbf{E2E Tests (Playwright):}
      \begin{itemize}
        \item Simulates complete user actions in real browsers.
        \item Best suited for verifying Leaflet map panning, zooming, and drag-and-drop file ingestion.
      \end{itemize}
  \end{itemize}
\end{frame}

% Slide 28: Vitest Configuration & Setup Mocks
\begin{frame}{Vitest Configuration \& Setup Mocks}
  \begin{itemize}
    \item \textbf{Headless DOM Simulation:} Vitest runs tests in `jsdom`. However, `jsdom` lacks definitions for several browser-native window features.
    \item \textbf{Mocking Setup File (`setup.ts`):}
      \begin{itemize}
        \item Mocked `window.localStorage` to support state persistence tests.
        \item Mocked `URL.createObjectURL` to prevent exceptions when generating upload image previews.
        \item Mocked Leaflet DOM resize behaviors that rely on physical window sizes.
      \end{itemize}
    \item \textbf{Benefit:} Tests run reliably in headless CLI environments (essential for CI/CD pipelines).
  \end{itemize>
\end{frame}

% Slide 29: API Layer Mocking: MSW (Mock Service Worker)
\begin{frame}[fragile]{API Layer Mocking: MSW Integration}
  \begin{itemize}
    \item \textbf{Strategy:} Instead of mocking global Axios instances (which couples tests to implementation details), we intercept network requests.
  \end{itemize}
\begin{lstlisting}[language=JavaScript]
const handlers = [
  http.get('http://localhost:3001/photos', () => {
    return HttpResponse.json([
      { id: 1, lat: 48.8566, lng: 2.3522, filename: 'paris.jpg', original_name: 'Paris Photo', ai_description: 'Eiffel Tower' }
    ]);
  })
];
const server = setupServer(...handlers);
\end{lstlisting}
  \begin{itemize}
    \item \textbf{Benefit:} Tests execute real Axios requests against a mock network layer, validating exact endpoint contracts.
  \end{itemize}
\end{frame}

% Slide 30: Code Auditing & Dependency Security
\begin{frame}{Code Auditing \& Dependency Security}
  \begin{itemize}
    \item \textbf{Vulnerability Management:} Run `npm audit` inside automated checks.
    \item \textbf{Resolving Critical Advisories:}
      \begin{itemize}
        \item Add `overrides` in `package.json` to force safe versions of sub-dependencies:
\begin{lstlisting}[language=JSON]
  "overrides": {
    "esbuild": "^0.25.0"
  }
\end{lstlisting}
        \item Upgraded development dependencies (e.g. `vite` to version `^6.4.2`) to resolve high-severity path traversal vulnerabilities.
      \end{itemize}
    \item \textbf{Benefit:} Protects the code supply chain from remote exploits.
  \end{itemize}
\end{frame}

% Slide 31: Dockerization & CI/CD Pipelines
\begin{frame}{Dockerization \& CI/CD Pipelines}
  \begin{itemize}
    \item \textbf{Containerization (Docker):}
      \begin{itemize}
        \item Package the Express backend and React frontend into standardized images.
        \item Prevents *"works on my machine"* issues by locking Node.js runtimes and package boundaries.
      \end{itemize}
    \item \textbf{CI/CD Automation (GitHub Actions):}
      \begin{itemize}
        \item Triggers on every Pull Request.
        \item Runs code audits (`npm audit`), checks formatting, executes unit tests (`vitest run`), and runs production build scripts.
      \end{itemize}
    \item \textbf{Outcome:} Guarantees code quality before merging to production branches.
  \end{itemize}
\end{frame}

% Slide 32: Production Scaling: SQLite Limitations
\begin{frame}{Production Scaling: SQLite Limitations}
  \begin{itemize}
    \item \textbf{Write Concurrency:} SQLite locks the entire database file on writes, causing write bottlenecks under high concurrent uploads.
    \item \textbf{Single-Process Dependency:} SQLite lives as a local file, making it impossible to scale the backend horizontally across multiple server nodes.
    \item \textbf{Missing Spatial Features:} SQLite lacks native, optimized spatial indices and geospatial operations required for location queries.
    \item \textbf{Production Path:} Migrate the database layer to PostgreSQL.
  \end{itemize}
\end{frame}

% Slide 33: Database Migration: PostgreSQL with PostGIS
\begin{frame}{PostgreSQL \& PostGIS Migration}
  \begin{itemize}
    \item \textbf{Why PostGIS:} PostGIS extends PostgreSQL with geographic objects and spatial query capabilities.
    \item \textbf{Spatial Data types:} Store coordinates inside a single `geom GEOGRAPHY(Point, 4326)` column.
    \item \textbf{SRID 4326:} Uses the standard WGS 84 spatial reference system (latitude/longitude coordinates used by GPS).
    \item \textbf{Automatic Updates:} Configure trigger functions to update the `geom` geography column automatically whenever `lat` or `lng` fields are updated.
  \end{itemize}
\end{frame}

% Slide 34: Spatial Indexing: GIST on Coordinate Points
\begin{frame}{Spatial Indexing: GIST on Coordinate Points}
  \begin{itemize}
    \item \textbf{The Problem:} Standard B-Tree indexes sort one-dimensional data. They cannot index two-dimensional geographic boundaries efficiently.
    \item \textbf{GIST Index (Generalized Search Tree):}
      \[ \text{CREATE INDEX idx\_photos\_geom ON photos USING GIST(geom);} \]
    \item \textbf{How it works:} GIST indexes group spatial shapes into bounding boxes. Queries can quickly traverse these boxes without executing full-table scans.
    \item \textbf{Performance Result:} Viewport queries execute in sub-millisecond times, even over millions of photo entries.
  \end{itemize}
\end{frame}

% Slide 35: Bounding Box Queries: SQL ST_MakeEnvelope
\begin{frame}[fragile]{Bounding Box Queries: SQL ST\_MakeEnvelope}
  \begin{itemize}
    \item \textbf{Algorithm:} Translate visible map coordinate boundaries directly to PostGIS envelopes.
    \item \textbf{SQL Spatial Query:}
\begin{lstlisting}[language=SQL]
SELECT id, filename, lat, lng, ai_description 
FROM photos 
WHERE geom && ST_MakeEnvelope(west, south, east, north, 4326);
\end{lstlisting}
    \item \textbf{Operator `\&\&`:} The bounding box intersection operator. Checks if the geographic coordinate point lies within the bounding box polygon.
    \item \textbf{Optimization:} Bypasses heavy geometry calculations to return viewport data instantly.
  \end{itemize}
\end{frame}

% Slide 36: Proximity Queries on Linear Assets
\begin{frame}[fragile]{Proximity Queries on Linear Assets}
  \begin{itemize}
    \item \textbf{Inspection Need:} Identify all inspection photos captured within 50 meters of a specific gas pipeline or power line segment.
    \item \textbf{PostGIS Proximity SQL:}
\begin{lstlisting}[language=SQL]
SELECT p.id, p.filename, p.original_name
FROM photos p, infrastructure_lines i
WHERE i.id = :lineId
  AND ST_DWithin(p.geom, i.geom, 50.0);
\end{lstlisting}
    \item \textbf{ST\_DWithin:} Computes the shortest distance between geometries on a geoid, utilizing spatial indexes for high-speed lookups.
    \item \textbf{Value:} Automates correlation of inspection photos to asset IDs.
  \end{itemize}
\end{frame}

% Slide 37: Cloud Storage Architecture: S3 & CDN Delivery
\begin{frame}{Cloud Storage Architecture: S3 \& CDN Delivery}
  \begin{itemize}
    \item \textbf{Storage Migration:} Move images out of the Express server folder to S3-compatible object storage (AWS S3, Cloudflare R2).
    \item \textbf{Presigned URLs:} Frontend requests upload $\to$ API returns a secure, short-lived presigned S3 URL $\to$ Client uploads image directly to S3.
    \item \textbf{Asset Delivery Pipeline:}
  \end{itemize}
  \begin{center}
    Client Request $\to$ CDN Edge Cache (Cloudflare/CloudFront) $\to$ S3 Bucket Source
  \end{center}
  \begin{itemize}
    \item \textbf{Benefit:} Reduces API server load, improves global image load times, and secures storage buckets behind access rules.
  \end{itemize}
\end{frame}

% Slide 38: Retrospective: JSDOM Map Mocking
\begin{frame}{Retrospective: JSDOM Map Mocking}
  \begin{itemize}
    \item \textbf{What was difficult:} Leaflet depends on actual browser DOM layout properties (e.g. `clientHeight`). Since JSDOM runs in a headless Node console, Leaflet maps threw rendering errors during unit tests.
    \item \textbf{How we solved it:} Created custom global mocks in `setup.ts` to bypass bounds calculations and mock Leaflet event registers.
    \item \textbf{Key Learning:} Map components are highly visual. While unit tests are great for forms and state logic, **Playwright End-to-End tests** are much better suited for validating interactive map views.
  \end{itemize}
\end{frame}

% Slide 39: Future Roadmap: WebGL Tiles & Edge Processing
\begin{frame}{Future Roadmap: WebGL Tiles \& Edge Processing}
  \begin{itemize}
    \item \textbf{WebGL Rendering (MapLibre GL):} Switch Leaflet DOM markers to WebGL canvas vector layers. Allows rendering $100,000+$ coordinate markers simultaneously without client lag.
    \item \textbf{Edge Computing Pipelines:} Move image resizing and AI processing to serverless edge workers (Cloudflare Workers/AWS Lambda). Offloads CPU-intensive image tasks from our core Express API.
    \item \textbf{Automated Computer Vision:} Integrate object detection models (YOLO, Segment Anything) into the upload stream to automatically detect asset defects (insulator damage, pipeline leaks) and tag them with coordinates.
  \end{itemize}
\end{frame}

% Slide 40: Summary & Open Q&A
\begin{frame}{Summary \& Open Q\&A}
  \begin{columns}
    \begin{column}{0.5\textwidth}
      \textbf{GeoPhoto Core Architecture:}
      \begin{itemize}
        \item Decoupled client-server design.
        \item Secure authentication & AI worker queues.
        \item Automated test suites (Vitest + MSW).
      \end{itemize}
      \textbf{Production Blueprint:}
      \begin{itemize}
        \item PostgreSQL + PostGIS spatial index query.
        \item S3 direct uploads & CDN caching.
      \end{itemize}
    \end{column}
    \begin{column}{0.5\textwidth}
      \begin{center}
        \Huge \textbf{Thank You!} \\
        \vspace{1cm}
        \large \textbf{Questions \& Answers}
      \end{center}
    \end{column}
  \end{columns}
\end{frame}

\end{document}
